\documentclass[UTF8]{ctexart}
\usepackage{amsmath}
\usepackage{booktabs}
\usepackage{geometry}
\usepackage{listings}
\usepackage{xcolor}
\usepackage{hyperref}
\usepackage{graphicx}

\geometry{a4paper, left=2.5cm, right=2.5cm, top=2.5cm, bottom=2.5cm}
\hypersetup{colorlinks=true, linkcolor=blue, urlcolor=blue, citecolor=blue}

\lstset{
  basicstyle=\ttfamily\small,
  breaklines=true,
  frame=single,
  keywordstyle=\color{blue},
  commentstyle=\color{gray},
  stringstyle=\color{teal},
  showstringspaces=false,
  tabsize=2,
  columns=flexible,
}

\title{\textbf{系统工具开发课程实验报告（第三周）}}
\author{廖世嘉 \quad 25020007073 \\ 25级计算机科学与技术}
\date{2026年9月7日}

\begin{document}
\maketitle
\tableofcontents
\newpage

\section{实验基本信息}

\begin{table}[htbp]
\centering
\caption{实验基本信息}
\label{tab:info}
\begin{tabular}{ll}
\toprule
项目 & 内容 \\
\midrule
姓名 & 廖世嘉 \\
学号 & 25020007073 \\
专业 & 25级计算机科学与技术 \\
实验环境 & Linux / Python / build / venv / PyTorch / pytest \\
实验日期 & 2026年9月7日 \\
GitHub仓库 & \href{https://github.com/MrLLLiao/SystemToolkitDevCourse}{https://github.com/MrLLLiao/SystemToolkitDevCourse} \\
\bottomrule
\end{tabular}
\end{table}

\section{实验目的}

本实验围绕 Python 包打包与分发、智能体编程、协作材料质量改写以及 PyTorch 训练循环修复四个主题展开。通过四个连续实验，掌握从源码构建 wheel 并在干净环境安装的流程，体验编程智能体的测试驱动修复循环，学习将低质量协作材料改写为可执行信息，以及补全 PyTorch 线性回归训练循环。

\section{实验环境与工具}

本实验在 Linux 和 Windows 环境中完成，主要使用 Python 的 build 模块、venv 虚拟环境、pip 安装、pytest 测试框架、ruff 代码检查工具以及 PyTorch CPU 版本。实验目录为 SystemToolkitDevCourse，四个实验分别位于 q09、q10、q11 和 q12。

\section{实验九：从源码构建并在干净环境安装 Wheel}

\subsection{实验要求}

在 q09 中创建一个最小 Python 命令行包 greetlab，包含 src/greetlab/\_\_init\_\_.py、cli.py 和 pyproject.toml。在课程环境中运行 python -m build 生成 wheel，新建一个干净虚拟环境只从生成的 wheel 安装，切换到 q09 之外运行 sdt-greet --name <学号> 确认输出正确。

\subsection{包结构}

项目目录结构如下：

\begin{verbatim}
q09/
├── pyproject.toml
└── src/
    └── greetlab/
        ├── __init__.py
        └── cli.py
\end{verbatim}

cli.py 实现了一个简单的命令行问候工具：

\begin{lstlisting}[language=Python]
import argparse

def main():
    p = argparse.ArgumentParser()
    p.add_argument("--name", required=True)
    a = p.parse_args()
    print(f"Hello, {a.name}!")
\end{lstlisting}

pyproject.toml 配置了构建系统和入口点：

\begin{lstlisting}
[build-system]
requires = ["setuptools>=68"]
build-backend = "setuptools.build_meta"

[project]
name = "greetlab-25020007073"
version = "0.1.0"

[project.scripts]
sdt-greet = "greetlab.cli:main"

[tool.setuptools.packages.find]
where = ["src"]
\end{lstlisting}

\subsection{构建与安装验证}

使用 python -m build --no-isolation 生成 wheel：

\begin{lstlisting}[language=bash]
cd q09
python -m build --no-isolation
# Successfully built greetlab_25020007073-0.1.0.tar.gz
# and greetlab_25020007073-0.1.0-py3-none-any.whl
\end{lstlisting}

创建干净虚拟环境并从 wheel 安装：

\begin{lstlisting}[language=bash]
python3 -m venv /tmp/q09_test_venv
source /tmp/q09_test_venv/bin/activate
pip install ./q09/dist/greetlab_25020007073-0.1.0-py3-none-any.whl
\end{lstlisting}

\subsection{结果}

在 q09 之外运行 sdt-greet --name 25020007073 输出：

\begin{lstlisting}
Hello, 25020007073!
\end{lstlisting}

确认 wheel 包安装正确，入口点命令 sdt-greet 可用。

\section{实验十：让编程智能体进入可验证的修复循环}

\subsection{实验要求}

复制 q09 为 q10，添加一个失败测试：当 name 只含空白字符时，main 应以 SystemExit(2) 结束。先运行测试确认失败，再向编程智能体说明目标、约束和测试命令，要求它修改实现并运行测试。人工检查 diff，撤销无关修改，并再次运行测试。在 ai\_log.md 中用不超过 5 行记录核心提示、智能体改动和人工验证。

\subsection{失败测试}

创建 test\_cli.py 包含正常姓名测试和空白姓名测试：

\begin{lstlisting}[language=Python]
import pytest
from greetlab.cli import main

def test_normal_name(capsys):
    import sys
    sys.argv = ["sdt-greet", "--name", "alice"]
    main()
    out = capsys.readouterr().out
    assert "Hello, alice!" in out

def test_blank_name_exits_nonzero():
    import sys
    sys.argv = ["sdt-greet", "--name", "   "]
    with pytest.raises(SystemExit) as exc_info:
        main()
    assert exc_info.value.code == 2
\end{lstlisting}

运行测试确认 test\_blank\_name\_exits\_nonzero 失败：

\begin{lstlisting}
FAILED test_blank_name_exits_nonzero - Failed: DID NOT RAISE SystemExit
1 failed, 1 passed
\end{lstlisting}

\subsection{智能体修复}

向编程智能体发出提示：修复 greetlab cli.py，当 --name 只含空白字符时 main 应以 SystemExit(2) 退出。约束：不改动参数解析逻辑，测试命令 PYTHONPATH=src pytest test\_cli.py。

智能体在 main() 中 parse\_args() 后添加了两行：

\begin{lstlisting}[language=Python]
if not a.name.strip():
    raise SystemExit(2)
\end{lstlisting}

\subsection{人工验证与记录}

检查 diff，确认仅涉及 cli.py 第 8--9 行，无无关修改。运行 pytest 两项测试均通过。ai\_log.md 记录如下：

\begin{lstlisting}
# AI Log
提示：修复 greetlab cli.py，当 --name 只含空白字符时 main 应以 SystemExit(2) 退出。
智能体改动：在 main() 中 parse_args() 后添加 if not a.name.strip(): raise SystemExit(2)，共增加2行。
人工验证：diff 仅涉及 cli.py 第8-9行，无无关修改。pytest 2 passed，空白姓名测试通过。
\end{lstlisting}

\section{实验十一：把"无法处理"的协作材料改成可执行信息}

\subsection{实验要求}

围绕"空姓名仍输出问候语"这一缺陷，改写三段质量较差的协作材料：Issue（Windows上运行不了，尽快修复）、提交信息（fix bug）、评审意见（这里写得不好，重写）。在 communication.md 中重写，包含环境、复现命令、期望结果和实际结果，未知信息写"待确认"。全文控制在 400 字以内。

\subsection{重写结果}

\textbf{Issue 重写}（包含环境、复现命令、期望/实际结果）：

\begin{itemize}
\item 标题：sdt-greet --name " " 输出非空问候语且退出码为0
\item 环境：Windows 11，Python 3.x，greetlab-25020007073 0.1.0
\item 复现命令：sdt-greet --name " "
\item 期望结果：当 name 只含空白字符时，程序应以非零退出码退出，不输出问候语
\item 实际结果：程序输出 Hello, ! 并以退出码0结束
\item 待确认：是否需要对 name 做 Unicode 空白字符（如全角空格）的处理
\end{itemize}

\textbf{提交信息重写}（祈使语气标题，正文说明问题和解决方案）：

\begin{lstlisting}
fix: reject whitespace-only name and exit with code 2

When --name contains only whitespace, the program previously printed
"Hello, !" and exited 0. Add validation in main() to detect blank
names, print an error message, and exit with SystemExit(2).
\end{lstlisting}

\textbf{评审意见重写}（标注 Blocking，指出具体行为、风险和建议动作）：

Blocking —— cli.py:main() 在 name 为纯空白时未做校验，导致输出无意义问候语且退出码为0，违反 CLI 约定。风险：下游脚本可能依赖退出码判断成功/失败，当前行为会掩盖错误。建议动作：在 parse\_args() 后增加 name.strip() 校验，若为空则 raise SystemExit(2)，并补充对应测试用例。

\subsection{结果}

重写后的三段材料均包含具体的环境信息、复现步骤、期望与实际结果对比，以及明确的修复建议，全文控制在 400 字以内，保持专业、简洁。

\section{实验十二：修复一个可复现的线性回归训练循环}

\subsection{实验要求}

将给定的 PyTorch 线性回归代码保存为 q12/train.py，补全两个 TODO：清空梯度、反向传播、更新参数；进入评估模式并在 no\_grad 中打印最终损失、weight 和 bias。最终损失应小于 0.001，不得直接把 weight 和 bias 赋值为 3 和 $-1$。

\subsection{原始代码与补全}

原始代码中训练循环缺少梯度清空、反向传播和参数更新步骤：

\begin{lstlisting}[language=Python]
for _ in range(200):
    pred = model(x)
    loss = loss_fn(pred, y)
    # TODO：清空梯度、反向传播、更新参数
    # TODO：进入评估模式并在no_grad中打印最终损失、weight和bias
\end{lstlisting}

补全后的完整训练循环：

\begin{lstlisting}[language=Python]
for _ in range(200):
    pred = model(x)
    loss = loss_fn(pred, y)
    opt.zero_grad()
    loss.backward()
    opt.step()

model.eval()
with torch.no_grad():
    final_pred = model(x)
    final_loss = loss_fn(final_pred, y)
    print(f"final_loss={final_loss.item():.6f}")
    print(f"weight={model.weight.item():.6f}")
    print(f"bias={model.bias.item():.6f}")
\end{lstlisting}

\subsection{结果}

运行 train.py 输出：

\begin{lstlisting}
final_loss=0.000000
weight=2.999997
bias=-1.000000
\end{lstlisting}

最终损失为 0.000000，小于 0.001 的要求。weight 接近 3，bias 接近 $-1$，与目标函数 $y = 3x - 1$ 一致。训练过程通过 SGD 优化器（学习率 0.1）在 200 轮迭代中收敛，未直接赋值参数。

\section{实验结果汇总}

\begin{table}[htbp]
\centering
\caption{四个实验的完成情况}
\label{tab:summary}
\begin{tabular}{cll}
\toprule
题号 & 实验主题 & 验证结果 \\
\midrule
9 & 打包与安装 Wheel & wheel 构建成功，干净环境安装并运行正确 \\
10 & 智能体修复循环 & 失败测试确认，智能体修复后 2 项测试通过 \\
11 & 协作材料改写 & Issue/提交/评审三段材料均改写为可执行信息 \\
12 & PyTorch 训练循环 & 补全 zero\_grad/backward/step，损失 $<0.001$ \\
\bottomrule
\end{tabular}
\end{table}

\section{问题与解决方法}

\subsection{build 模块的 venv 依赖问题}

在 Windows Python 中运行 python -m build 时，由于缺少 venv 模块导致构建失败。解决方法是使用 --no-isolation 参数跳过隔离环境创建，直接使用当前 Python 环境中的 setuptools 进行构建。

\subsection{PyTorch 安装问题}

WSL 环境中的 Python 3.14 没有预装 PyTorch，且 externally-managed-environment 限制了全局安装。解决方法是在 Windows Python 中通过 CPU 专用索引安装 PyTorch（pip install torch --index-url https://download.pytorch.org/whl/cpu），安装后成功运行训练脚本。

\subsection{智能体修改的范围控制}

在实验十中，需要确保智能体的修改仅限于必要范围。通过人工检查 diff，确认修改仅涉及 cli.py 中新增的两行校验代码，没有对参数解析逻辑或其他文件产生无关变更。

\section{实验总结}

通过本次实验，我完整体验了从源码打包到 wheel 分发的流程，理解了 pyproject.toml 中构建系统、入口点和包发现的配置方式。实验十让我认识到编程智能体在测试驱动修复中的价值：只要提供明确的目标、约束和测试命令，智能体可以完成最小化修复，但仍需人工审查 diff 以确保不引入无关变更。

实验十一强化了协作材料的质量意识：Issue 应包含可复现的环境和命令，提交信息应使用祈使语气并说明"为什么"，评审意见应标注严重级别并给出具体建议。实验十二则让我实践了 PyTorch 训练循环的标准流程：zero\_grad $\to$ backward $\to$ step 的三步序列，以及评估模式下使用 no\_grad 防止梯度计算的正确做法。

本周的核心认识是：无论是包安装、智能体修复还是模型训练，"可验证性"都是关键——wheel 需要在干净环境中验证安装，修复需要通过测试验证，训练需要通过损失值验证收敛。

\end{document}
